\documentclass[12pt,a4paper]{article}
\usepackage[UTF8,heading=true]{ctex}
\usepackage{geometry}
\usepackage{amsmath,amssymb}
\usepackage{enumitem,hyperref,fancyhdr,setspace}
\usepackage{ragged2e}

\geometry{left=2.5cm,right=2.5cm,top=2.5cm,bottom=2.5cm}
\setlength{\headheight}{15pt}
\setstretch{1.5}
\setlength{\emergencystretch}{3em}
\hfuzz=20pt
\sloppy
\hypersetup{hidelinks}

\pagestyle{fancy}
\fancyhf{}
\fancyhead[C]{发明专利权利要求书}
\fancyfoot[C]{\thepage}
\renewcommand{\headrulewidth}{0pt}

\begin{document}

\begin{center}
    {\Large\heiti 中国移动专利申请} \\[0.5em]
    {\Huge\heiti 权利要求书} \\[0.8em]
\end{center}

\noindent\textbf{发明名称：}量子-经典异构光纤网络的动态波分复用与路由控制方法

\vspace{0.8em}

\begin{enumerate}[
    label=\arabic*.,
    leftmargin=0pt,
    itemindent=2em,
    labelsep=0.5em,
    align=left,
    itemsep=0.6em,
    topsep=0pt
]

\item 一种面向量子-经典异构光纤网络的动态波分复用与路由控制方法，其特征在于，应用于量子控制器/编排器，所述异构光纤网络在同一光纤中承载量子信号与经典波分复用业务信号，并与 ROADM/OTN 设备协同；所述方法包括：
（1）采集链路段级的经典信道功率谱、波长占用信息、放大器状态与告警信息，以及量子链路测量指标，所述量子链路测量指标至少包括量子误码率 QBER、单光子计数率、暗计数率、探测器饱和度和密钥率；
（2）基于所述采集数据构建链路段-波长-时间窗三级噪声感知数据，形成三维噪声张量，并计算候选路径、候选量子波长以及候选发射时间窗在给定探测门控参数与纠错配置下的预测 QBER 与预测安全密钥率；
（3）在满足量子业务质量约束与现网扰动约束的条件下，对路由 $p$、量子波长 $\lambda$、发射时间窗 $t$、纠错配置 $e$ 和探测门控参数 $g$ 进行联合优化，得到目标决策 $(p^\star,\lambda^\star,t^\star,e^\star,g^\star)$；
（4）按先轻后重的分层执行策略下发重配置：优先执行协议层轻量动作以切换发射时间窗和/或探测门控参数和/或纠错配置；当仍不满足所述量子业务质量约束时，再执行量子波长迁移和/或路径迁移；
（5）在执行迁移前生成上一稳定配置快照，并在迁移窗口内保持新旧配置的重叠保护；当重配置失败或重配置后量子业务质量未达标时，在保护窗口内执行回滚并记录审计信息。

\item 如权利要求1所述的方法，其特征在于，所述发射时间窗 $t$ 为发射侧与接收侧一致的离散时间区间 $[t_k, t_k+\Delta t]$，并在每个时间窗内统计经典业务占用与功率谱密度以形成占用张量 $O(p,\lambda,t)$，进而基于所述占用张量估计噪声光子数并预测 $QBER(p,\lambda,t,g,e)$ 与安全密钥率 $R_{\mathrm{sec}}(p,\lambda,t,g,e)$。

\item 如权利要求1所述的方法，其特征在于，所述构建三级噪声感知数据包括对拉曼散射噪声、放大自发辐射噪声、串扰与泄露以及链路预算裕量进行联合建模，并生成用于候选决策筛选的综合风险指标。

\item 如权利要求1所述的方法，其特征在于，所述联合优化以效用函数为目标函数，所述效用函数至少同时表征预测安全密钥率、预测 QBER、网络重构代价以及对现网经典业务的扰动成本。

\item 如权利要求1所述的方法，其特征在于，所述量子业务质量约束至少包括 $QBER \leq QBER_{\max}$ 且 $R_{\mathrm{sec}} \geq R_{\min}$，所述现网扰动约束至少包括任一次波长迁移、路径迁移或门控参数调整导致的扰动代价满足 $P_{\mathrm{impact}} \leq P_{\max}$。

\item 如权利要求1所述的方法，其特征在于，所述分层执行策略至少包括：
L1 协议层轻量动作：调整发射时间窗、探测门控宽度与门控相位、诱骗态比例以及纠错码率与块长；
L2 光层中量动作：在不改变端到端拓扑的前提下执行量子波长迁移，并与滤波/保护带策略及协议层参数同步切换；
L3 路由级重构动作：执行路径迁移，必要时联合量子波长迁移。

\item 如权利要求6所述的方法，其特征在于，仅当预测劣化趋势持续 $K$ 个采样周期且效用增益超过阈值时，才允许从 L1 升级至 L2 或 L3；并且在波长迁移或路径迁移后至少驻留预设最小驻留时间或预设采样周期数后才允许再次触发同级光层重构。

\item 如权利要求1所述的方法，其特征在于，在遥测缺失或数据异常的情况下进入保守模式，所述保守模式包括冻结光层重构并仅允许执行 L1 协议层轻量动作，同时提高切换触发门限并触发二次采样或人工复核流程。

\item 如权利要求1所述的方法，其特征在于，当预测模型判定为短时噪声尖峰且持续时间小于阈值时，不触发 L2 或 L3 动作，而是通过时间窗错避、门控收窄与纠错增强进行快速抑制以避免频繁迁移。

\item 如权利要求1所述的方法，其特征在于，当检测到探测器计数率或饱和度达到阈值时，优先降低发射占空比和/或缩短探测门控宽度和/或调整诱骗态配置，并在必要时切换至更低噪声的量子波长或路径以降低后脉冲风险。

\item 如权利要求1所述的方法，其特征在于，所述重叠保护包括在保护窗口内保持旧配置与新配置的短时并行可用，并在收到 ROADM/WSS/OTN 的执行回执与量子测量指标满足门限后再释放旧配置。

\item 一种面向量子-经典异构光纤网络的动态波分复用与路由控制系统，其特征在于，包括：
经典噪声遥测模块，用于采集 ROADM/OTN 的功率谱、端口功率、波长占用、放大器状态与告警信息；
量子链路测量模块，用于采集 QKD 收发端的 QBER、计数率、暗计数率、探测门控状态、饱和度和密钥率；
噪声建模模块，用于构建链路段-波长-时间窗三级噪声张量并输出预测指标；
联合优化引擎，用于对路由 $p$、量子波长 $\lambda$、发射时间窗 $t$、纠错配置 $e$ 与探测门控参数 $g$ 进行联合求解并输出重构计划；
重配置编排器，用于按分层执行策略向 ROADM/WSS/OTN 设备及 QKD 设备下发重配置指令，并执行重叠保护与回滚；
策略审计模块，用于记录调度依据、阈值、影响评估与执行结果。

\item 如权利要求12所述的系统，其特征在于，所述经典噪声遥测模块通过 NETCONF、gNMI、SNMP 或网管北向接口获取所述功率谱与占用信息，所述重配置编排器通过 SDN 控制器接口或网管接口下发对 ROADM/WSS/OTN 的重构指令，且所述量子链路测量模块通过量子设备控制/遥测接口获取所述量子链路测量指标。

\item 如权利要求12所述的系统，其特征在于，所述系统向量子密钥管理系统或业务编排系统输出可用密钥率、业务降级状态和迁移计划，以形成可测、可算、可控、可审计的闭环。

\item 如权利要求12所述的系统，其特征在于，所述重配置编排器在执行任一次波长迁移或路径迁移前生成可回滚的上一稳定配置快照，并在预设保护窗口内完成重构；当执行回执异常或量子业务质量在预设采样周期内未改善时，触发回滚并将对应候选动作标记为短期禁用。

\item 一种电子设备，包括存储器、处理器及存储在所述存储器中并可在所述处理器上运行的计算机程序，其特征在于，所述处理器执行所述计算机程序时实现如权利要求1至11任一项所述的方法。

\item 一种计算机可读存储介质，其上存储有计算机程序，其特征在于，所述计算机程序被处理器执行时实现如权利要求1至11任一项所述的方法。

\end{enumerate}

\end{document}

